\subsubsection{Backend -- Spring Boot}
\label{subsubsec:backend-spring-boot}

Nhóm lựa chọn Spring Boot để xây dựng backend của hệ thống chat. Spring Boot
hỗ trợ phát triển ứng dụng dựa trên Spring thông qua cơ chế tự động cấu hình
và quản lý các thành phần của ứng dụng, giúp giảm phần cấu hình thủ công
khi khởi tạo dự án.

Trong đồ án, Spring Boot được định hướng sử dụng như sau:

\begingroup
\small
\begin{longtable}{|L{4.6cm}|L{\dimexpr\textwidth-4.6cm-4\tabcolsep-3\arrayrulewidth\relax}|}
  \caption{Định hướng áp dụng Spring Boot trong đồ án}
  \label{tab:ap-dung-spring-boot}\\
  \hline
  \textbf{Đặc điểm} & \textbf{Cách áp dụng trong đồ án}\\
  \hline
  \endfirsthead
  \hline
  \textbf{Đặc điểm} & \textbf{Cách áp dụng trong đồ án}\\
  \hline
  \endhead
  Quản lý thành phần và Dependency Injection &
  Quản lý các đối tượng xử lý nghiệp vụ, truy cập dữ liệu và các thành phần
  phục vụ server; cung cấp chúng cho những nơi cần sử dụng\\
  \hline
  Tự động cấu hình &
  Hỗ trợ cấu hình các thành phần dựa trên thư viện và thông tin được khai báo
  trong dự án\\
  \hline
  Cấu hình bên ngoài &
  Khai báo cổng TCP và thông tin kết nối PostgreSQL trong file cấu hình,
  hạn chế ghi trực tiếp vào mã nguồn\\
  \hline
  Tổ chức mã nguồn &
  Kết hợp với cách phân chia service (xử lý nghiệp vụ) và DAO
  (Data Access Object, lớp truy cập dữ liệu) để tách xử lý nghiệp vụ
  khỏi truy vấn dữ liệu\\
  \hline
  Đóng gói ứng dụng &
  Hỗ trợ đóng gói backend thành ứng dụng Java có thể chạy độc lập\\
  \hline
\end{longtable}
\endgroup

Spring Boot không bắt buộc ứng dụng phải hoạt động như một web server.
Với yêu cầu của đồ án, backend được định hướng chạy dưới dạng ứng dụng
không có web server, kết hợp \texttt{ServerSocket} và \texttt{Socket} của Java
để giao tiếp TCP với các client. Spring Boot quản lý các thành phần của backend,
còn socket đảm nhiệm việc truyền nhận dữ liệu.

Trong Version 1, nhóm xác định vai trò và cách áp dụng Spring Boot ở mức
thiết kế. Cấu hình, mã nguồn backend và các cơ chế xử lý cụ thể sẽ được
triển khai trong phiên bản tiếp theo.
